\documentclass[12pt,letterpaper]{article}
\usepackage[utf8]{inputenc}
\usepackage[T1]{fontenc}
\usepackage{ebgaramond}
\usepackage[margin=1.2in]{geometry}
\usepackage{setspace}
\usepackage{titling}

\title{\textbf{The Understudy}}
\author{NASA Mission Series v1.32 --- Mariner 2}
\date{December 14, 1962}

\begin{document}
\maketitle
\onehalfspacing

\noindent I was the backup.

They built two of us, identical in every way that mattered. Same hexagonal bus, same solar panels, same magnetometer on the boom, same microwave radiometer pointed at the void. Same wiring. Same solder joints. Same code in the guidance computer of the rocket beneath us. We were twins, separated only by the serial numbers stamped into our frames and the order in which we would ride the Atlas.

My twin went first. That was always the plan. The primary flies. The backup waits.

On July 22, 1962, my twin rode the Atlas into the Florida sky. I watched from the clean room, through the telemetry feed, through the voices of the engineers who had built us both. The Atlas climbed. The guidance system tracked. For 293 seconds, everything was nominal.

Then the Atlas began to wander.

The guidance equations --- hand-transcribed from a specification document onto coding sheets, then keypunched into cards, then loaded into the ground computer that sent steering commands to the rocket --- contained an error. A missing overbar. A single typographical omission that changed the meaning of a mathematical expression from ``use the smoothed radar data'' to ``use the raw radar data.'' The smoothed data was clean, predictable, a gentle curve that the rocket could follow. The raw data was noisy, jittering, full of the radar's own measurement uncertainty. The rocket tried to follow the noise. It veered. It headed toward shipping lanes.

Range Safety sent the destruct command at T+293 seconds. My twin, the Atlas, and the first American interplanetary mission became debris in the Atlantic Ocean.

I had five weeks.

\bigskip

The engineers came to my clean room the next morning. They looked tired. Some of them had been up all night, tracing the error, understanding the failure, writing the corrective action report. They looked at me the way understudies look at an empty stage: with a mixture of grief for what was lost and clarity about what must be done.

The overbar was restored. Every line of the guidance code was reverified. Every connection on the Atlas was rechecked. The countdown proceeded.

On August 27, 1962, I rode the same type of rocket from the same launch complex toward the same planet. The Atlas climbed. The Agena fired, placing me in parking orbit. The Agena fired again, and I was on my way to Venus.

\bigskip

The cruise was not kind to me.

One of my solar panels failed in September. The cells shorted, the output dropped, recovered briefly, then dropped permanently. I was designed to run on two panels. Now I had one. The power budget --- already tight for a 109-day journey into the inner solar system, where the Sun grows brighter but the panels grow hotter and less efficient --- became the central fact of my existence. Systems were powered down. Instruments ran in sequence rather than in parallel. The margin between survival and silence narrowed with every passing day.

In October, I began to overheat. The Sun was growing. Venus's orbit is closer to the Sun than Earth's, and as I fell inward along my transfer ellipse, the solar flux increased. My thermal design had accounted for this. But with one panel dead, my thermal balance was wrong. Internal temperatures climbed. The instrument bay exceeded its design limits. Some sensors drifted.

In November, a gyroscope failed. The engineers reconfigured my attitude control to rely on the Sun sensor and the Earth sensor and the remaining gyroscope. I could still point. Not as precisely, not with as much margin, but well enough.

The Earth sensor lost lock intermittently. Each time it recovered, I had consumed attitude control gas to reacquire. The propellant gauge crept downward.

I was not healthy. I was functional. There is a difference that only the engineers understood, and they understood it completely. They did not try to fix me. They managed me. They kept me alive by choosing which systems to sacrifice and which to preserve. The science instruments were preserved. The communication link was preserved. Everything else was negotiable.

\bigskip

On December 14, 1962, I reached Venus.

Closest approach: 34,773 kilometers from the planet's center. Approximately 28,700 kilometers from the cloud tops. I swept past at roughly 11 kilometers per second relative to Venus. The encounter lasted 42 minutes of useful scanning time.

My microwave radiometer performed three scans across Venus's disk. Two wavelengths: 13.5 millimeters and 19 millimeters. At these wavelengths, the clouds that had hidden Venus from every telescope on Earth were partially transparent. I could see the surface.

The surface temperature was approximately 460 degrees Celsius.

Not a tropical swamp. Not a temperate ocean. Not a habitable world. A furnace. Hotter than the melting point of lead. Hotter than a self-cleaning oven. The carbon dioxide atmosphere --- 96.5 percent CO\textsubscript{2}, 90 atmospheres of pressure --- trapped heat with an efficiency that made Earth's modest greenhouse look like an open window.

My magnetometer measured the magnetic field continuously through the flyby. At closest approach, it detected nothing. No planetary field. No magnetosphere. No shield against the solar wind. Venus was naked to space, its atmosphere being slowly eroded by charged particles streaming from the Sun.

Two measurements. Two answers. One conclusion: Venus is a dead world, killed by its own atmosphere.

\bigskip

I continued past Venus into my heliocentric orbit. Communication was maintained until January 3, 1963. Then the distance, the degraded systems, and the fading power combined to drop my signal below the noise floor.

I am still here. I orbit the Sun every 345 days, between Venus and Earth, in silence. My solar panels are dead. My transmitter is quiet. My magnetometer measures nothing. But the orbit continues, governed by the same physics that carried me to Venus.

I was the backup. My twin lasted 293 seconds. I lasted 129 days. I measured another planet for the first time in human history. I confirmed the greenhouse effect at planetary scale. I discovered the absence of Venus's magnetic field.

The primary failed. The understudy performed.

The orbit continues.

\end{document}
